\section{Local inverses and Gaussian normalization}
\label{sec:local-covariance}
The Gaussian covariance used in the expansion is assembled from local
principal inverses. It is not the full inverse of the precision matrix.
We prove its positivity, estimate the resulting error, and express the
normalization through local terms and convergent correction series.
This follows the real interpolation from the identity used in the
construction.

Let $I$ be a finite coordinate set, with anchors $a_i$ on a coarse
four-dimensional torus. Assume at most $m$ coordinates have any
given anchor. For a real matrix $T$ define
\[
 \|T\|_\eta=\max\left\{
 \sup_i\sum_j e^{\eta\operatorname{dist}(a_i,a_j)}|T_{ij}|,
 \sup_j\sum_i e^{\eta\operatorname{dist}(a_i,a_j)}|T_{ij}|
 \right\}.
\]
On an empty coordinate set both suprema are zero. The triangle
inequality for distance and an interchange of finite sums give
$\|ST\|_\eta\le\|S\|_\eta\|T\|_\eta$ for $\eta\ge0$.
Transposition preserves this norm, and decreasing $\eta$ cannot
increase it. It dominates the Euclidean operator norm by the
row-and-column Cauchy--Schwarz estimate used above.

Let $A=A^*$ and suppose
\begin{equation}\label{eq:covariance-data}
 A\ge r1,\qquad \|A\|_\eta\le Q,
 \qquad 0<r\le1,\quad Q\ge1,\quad\eta>0.
\end{equation}
Section~\ref{sec:sector-gram} supplies these constants for the
actual completed sector precision, after fixing $L$ and then
taking the large parameter sufficiently large. Every principal
submatrix has the same lower bound and weighted estimate: extend
a test vector by zero for the first assertion, and restrict the
nonnegative row and column sums for the second.

\subsection{Exponential decay of the inverse}
Set
\[
 \theta=\min\{1,r/(2Q)\},\qquad \gamma=\theta\eta,
 \qquad c_0=\min\{1,r/2\}.
\]
For every $u\ge0$, the matrix $A+u1$ is invertible and
\begin{equation}\label{eq:covariance-entry}
 e^{\gamma\operatorname{dist}(a_i,a_j)}
       |(A+u1)^{-1}_{ij}|\le\{c_0(1+u)\}^{-1}.
\end{equation}
We give the weighted argument, which is uniform in the number of
coordinates. Fix $j$ and put
$s_i=e^{\gamma\operatorname{dist}(a_i,a_j)}$.
For $0\le\theta\le1$ and $|v|\le\theta D$ with $D\ge0$,
convexity of the exponential gives
\[
 |e^v-1|\le\theta e^D.
\]
For positive $v$, use
$e^{\theta D}\le1+\theta(e^D-1)$; for negative $v$, use
$1-e^v\le e^{-v}-1$. The reverse triangle inequality for
distances now implies
\[
 |s_i/s_k-1|\le\theta e^{\eta\operatorname{dist}(a_i,a_k)}.
\]
Conjugate $A+u1$ by the positive diagonal matrix with entries
$s_i$. Its difference $F$ from $A+u1$ has entries
$(A+u1)_{ik}(s_i/s_k-1)$. The diagonal contribution is zero,
so both absolute row and column sums of $F$ are at most
$\theta Q\le r/2$. Thus
$|\langle w,Fw\rangle|\le(r/2)\|w\|_2^2$.

Let $y=(A+u1)^{-1}e_j$ and $w_i=s_i y_i$. The conjugated
equation is $(A+u1+F)w=e_j$, since $s_j=1$.
It follows that
\[
 (r/2+u)\|w\|_2^2\le w_j
 \le\|w\|_2.
\]
If $w\ne0$, division gives
$\|w\|_2\le(r/2+u)^{-1}\le\{c_0(1+u)\}^{-1}$;
the case $w=0$ is immediate. Each $|w_i|$ is bounded by this
norm, proving \eqref{eq:covariance-entry}. Positive definiteness
gave invertibility before this computation. The same proof applies
to every principal submatrix.

Put $k=\gamma/3$ and
\[
 S=c_0^{-1}mP_0(\gamma-2k),\qquad
 P_0(t)=\sum_{n\ge0}(2n+1)^4e^{-tn}.
\]
Summing \eqref{eq:covariance-entry} at $u=0$ with weight
$e^{2k\operatorname{dist}}$ gives $\|A^{-1}\|_{2k}\le S$.
The summation uses the anchor multiplicity $m$ and the same
distance-shell count as in Section~\ref{sec:sector-gram}.

\subsection{Principal windows and their error}
For each coordinate $j$, choose a window $W_j$ containing every
coordinate within distance $B$ of $a_j$ and no coordinate farther
than $B+2$. Define $P$ by its columns:
\[
 P_{ij}=\begin{cases}
 (A_{W_j}^{-1})_{ij},&i\in W_j,\\
 0,&i\notin W_j.
 \end{cases}
\]
Indices inside a window mean the corresponding principal coordinates.
All these inverses exist by \eqref{eq:covariance-data}. Applying
\eqref{eq:covariance-entry} separately to each column yields the
same pointwise exponential bound for $P$, hence
$\|P\|_{2k}\le S$. This also controls column sums even though
the windows vary with the column.

Let $E=AP-1$. If $i\in W_j$, the principal inverse equation
gives $E_{ij}=0$: the column $P_{\cdot j}$ is zero outside
$W_j$, so the full product at $i$ is exactly the principal product.
In particular $E_{ij}=0$ when the anchor distance is at most $B$.
The product bound gives $\|E\|_{2k}\le QS+1$.
On its remaining entries the weight can be reduced from $2k$ to
$k$ at cost $e^{-kB}$. Therefore
\begin{equation}\label{eq:window-error}
 \|E\|_k\le\epsilon_B:=e^{-kB}(QS+1),\qquad
 \|P-A^{-1}\|_k\le S\epsilon_B.
\end{equation}
The second inequality uses the exact identity
$P-A^{-1}=A^{-1}E$. Notice the order of the factors.

The physical covariance is the symmetric matrix
\begin{equation}\label{eq:local-covariance}
 C=\tfrac12(P+P^*).
\end{equation}
Symmetry of $A^{-1}$ and \eqref{eq:window-error} give
\[
 \|C\|_k\le S,\qquad
 \|C-A^{-1}\|_k\le\delta_B:=S\epsilon_B,
 \qquad \|AC-1\|_k\le Q\delta_B.
\]
The upper bound $A\le Q1$ implies $A^{-1}\ge Q^{-1}1$.
It can also be checked directly: for $w=A^{-1}v$ and
$x=Q^{-1}v$, positivity of
$\langle w-x,A(w-x)\rangle$ and
$\langle x,Ax\rangle\le Q\|x\|^2$ give
$\langle v,A^{-1}v\rangle\ge Q^{-1}\|v\|^2$.
Consequently, once $\delta_B\le1/(2Q)$,
\begin{equation}\label{eq:local-covariance-floor}
 C\ge(2Q)^{-1}1.
\end{equation}
The actual bond windows are colour independent: a coordinate belongs
to $W_j$ when some coarse endpoint of its bond is within $B$ of
some endpoint of the bond of $j$. Each endpoint is within one
of its bond's parent anchor. This proves the required $B$ and
$B+2$ bounds and shows that complete colour blocks occur together.
The matrix $C$ has range at most $B+2$. Its closeness to $A^{-1}$
is a conclusion of the preceding estimates, not a replacement of
one matrix by the other.

For the sector precision, inactive rows and columns of $A$ are
those of the identity. This property passes to every principal
inverse, hence to $P$ and $C$. Also $M^*b$ vanishes on inactive
coordinates, so the chosen local mean $-P^*M^*b$ vanishes there.
These observations justify restricting the local Gaussian to the
active coordinates without changing its physical entries.

\subsection{Ordered inverse series}
If a matrix $E$ satisfies $\|E\|_k\le\epsilon<1$, then
\[
 (1+E)^{-1}=\sum_{n=0}^\infty(-E)^n.
\]
Indeed the partial geometric identity has remainder $(-E)^N$,
whose norm tends to zero; the series converges absolutely in the
finite matrix space. For $E=AP-1$ this gives
\begin{equation}\label{eq:local-inverse-series}
 A^{-1}=\sum_{n\ge0}P(-E)^n.
\end{equation}
For $D=AC-1$ the corresponding identity is
\begin{equation}\label{eq:covariance-inverse-series}
 C^{-1}=\sum_{n\ge0}(-D)^nA,
 \qquad C^{-1}-A=\sum_{n\ge0}(-D)^{n+1}A.
\end{equation}
Multiply these formulas by $A$ on the left and by $C$ on the
right, respectively, to verify them. No commutation of $P$, $C$
and $A$ is asserted. Since $B\asymp(\log x)^2$ and $k>0$ is
fixed after $L$, $\epsilon_B$ and $\delta_B$ have every negative
upper exponent. Both convergence conditions can thus be met on
one sufficiently late tail, uniformly over sectors and periods.

\subsection{Logarithms as local terms and correction series}
For $0\le s\le1$ put $A_s=1+s(A-1)$. The hypotheses
\eqref{eq:covariance-data} hold for $A_s$ with the same
$r,Q,\eta$, since $r\le1\le Q$. Let $P_s$ be its window
inverse and $E_s=A_sP_s-1$. All bounds above are uniform in $s$.
Define the local diagonal terms and their correction terms by
\[
 \ell_i=\int_0^1[(A-1)P_s]_{ii}\,ds,
 \qquad
 \ell_{i,n}=\int_0^1[(A-1)P_s(-E_s)^n]_{ii}\,ds.
\]
Thus $\ell_i=\ell_{i,0}$. We have the exact identity
\begin{equation}\label{eq:local-logdet}
 \log\det A=\sum_i\ell_i+
                   \sum_i\sum_{n\ge1}\ell_{i,n}.
\end{equation}
To justify it, the determinant's linear term at the identity is
the trace, as follows by expanding the finite determinant by
permutations. Factoring
\[
 A_{s+t}=(1+t(A-1)A_s^{-1})A_s
\]
therefore gives
$\frac{d}{ds}\log\det A_s=
\operatorname{tr}((A-1)A_s^{-1})$.
All $A_s$ are positive definite. Their inverses and principal
inverses vary continuously, by the adjugate formula and the
nonvanishing determinant. Integrate the derivative from zero to
one and substitute \eqref{eq:local-inverse-series}. Its diagonal
terms are bounded, uniformly in $s$, by
$(Q+1)S\epsilon_B^n$. This summable geometric bound justifies
termwise integration for each $i$. The coordinate sum is finite.
In particular,
\[
 |\ell_i|\le(Q+1)S,\qquad
 \sum_{n\ge1}|\ell_{i,n}|
 \le\frac{(Q+1)S\epsilon_B}{1-\epsilon_B}.
\]
These are per-coordinate bounds. Summing them over the entire
torus introduces its coordinate count; it does not yield a
per-seed estimate without the further support assignment.

There is a second determinant correction for the covariance $C$.
With $D=AC-1$ and $\|D\|_k<1$, the path $1+sD$ is invertible
for $s\in[0,1]$ by its Neumann series. Its real determinant
starts at one and never vanishes, so it remains positive.
The same determinant derivative and geometric domination give
\begin{equation}\label{eq:covariance-logdet}
 \log\det A+\log\det C
 =\sum_i\sum_{n\ge0}\int_0^1[D(-sD)^n]_{ii}\,ds.
\end{equation}
Here $1+D=AC$ and both $A,C$ have positive determinant, so
multiplicativity of the determinant justifies the sum of logarithms.
Although $D$ need not be symmetric, no spectral logarithm of $D$
is used. The integral of its $n$th term is
$(-1)^n(D^{n+1})_{ii}/(n+1)$, with absolute value at most
$\|D\|_k^{n+1}/(n+1)$.

\subsection{Local comparison of normalization terms}
For a physical bond $v$, group its colour coordinates and set
$J_v(A)=\sum_{i:\,i\text{ lies over }v}\ell_i$.
The term $\ell_i$ reads only the principal window $W_i$:
\[
 [(A-1)P_s]_{ii}
 =[(A_{W_i}-1)(1+s(A_{W_i}-1))^{-1}]_{ii}.
\]
This follows because column $i$ of $P_s$ vanishes outside $W_i$.
Suppose $A$ and $Z$ both satisfy \eqref{eq:covariance-data},
use the same windows, and differ entrywise by at most $e$ on
$W_i\times W_i$. The window has at most
$\Delta=m(2B+5)^4$ coordinates. Thus the unweighted row and
column norms of $A_{W_i}-Z_{W_i}$ are at most $\Delta e$.
The same bound holds for the difference of their interpolants.
The inverse identity, with its order retained, is
\[
 (Z_s)_{W_i}^{-1}-(A_s)_{W_i}^{-1}
 =(Z_s)_{W_i}^{-1}
       ((A_s)_{W_i}-(Z_s)_{W_i})(A_s)_{W_i}^{-1}.
\]
Each inverse has row-and-column norm at most $S$. Splitting
the difference of the two local products into a difference of
first factors and a difference of inverses bounds its norm by
\[
 \{S+(Q+1)S^2\}\Delta e.
\]
Integration and the colour sum give
\begin{equation}\label{eq:local-log-comparison}
 |J_v(A)-J_v(Z)|
 \le q_{\rm col}\{S+(Q+1)S^2\}\Delta e,
\end{equation}
provided the entrywise hypothesis holds on the window for each
colour over $v$.

Finally $J_v$ is invariant under bondwise orthogonal changes of
colour frame. The windows contain whole colour blocks, so such
a change restricts to an orthogonal matrix on every window and
commutes with the projection onto the colours over $v$.
Principal interpolants and their inverses are conjugated by this
matrix. The sum of diagonal entries over $v$ is the trace after
that projection; cyclicity of the finite trace proves invariance
before integration. This is why one compares the whole colour
trace, rather than assigning independent gauge-invariant meanings
to its individual diagonal terms.
